\section{Conclusion}

This thesis examined whether cryptocurrency information, and Bitcoin in particular, can help forecast changing dependency structures with conventional assets and provide a practically relevant warning signal for investors. To answer this question, a reproducible two-layer framework was developed. The first layer focused on forecasting time-varying rolling dependency between Bitcoin and a set of conventional assets. The second layer translated the output of the first stage into a stress-warning classification problem.

The empirical results show that the dependency process is forecastable. Across the experiment set, out-of-sample forecasting accuracy is consistently strong, especially for cross-crypto and medium-to-longer-window dependencies. However, the primary source of predictive power is the persistence of the dependency series itself. This explains why AR1 and naive persistence models dominate the average ranking. Rather than weakening the thesis, this finding clarifies the structure of the problem and provides a more honest interpretation of what the data actually support.

Machine-learning models still play an important role in the thesis. They often outperform the DCC-GARCH benchmark and provide a richer comparison framework, even when they do not beat the strongest persistence baselines on average. The investor signal layer further extends the contribution by showing that crypto-derived information can be converted into an early warning signal for stress regimes in conventional markets. The signal is moderate rather than dramatic, but it is meaningful enough to justify interpretation as a risk-management overlay.

The thesis therefore makes a defensible academic contribution and a limited but coherent practical contribution. Academically, it shows that time-varying intermarket dependency between Bitcoin and conventional assets can be forecast in a disciplined out-of-sample framework. Practically, it shows that such forecasts can inform an investor-facing warning mechanism, especially for broad equity markets.

Future work should extend the horizon set, enrich the feature space with macro and on-chain variables, test alternative dependence targets, and integrate the signal layer into a portfolio-level backtesting framework. Even in its current form, however, the thesis provides a meaningful answer to the research problem and demonstrates that cryptocurrency data can contribute to understanding, forecasting, and managing changing cross-market risk.
